\documentclass[11pt,a4paper]{article}
\usepackage[utf8]{inputenc}
\usepackage[russian]{babel}
\usepackage{geometry}
\usepackage{booktabs}
\usepackage{array}
\usepackage{graphicx}
\usepackage{amsmath}
\usepackage{hyperref}
\usepackage{enumitem}
\usepackage{float}

\geometry{top=2cm, bottom=2cm, left=2.5cm, right=2.5cm}
\setlength{\parindent}{0pt}
\setlength{\parskip}{0.5em}

\title{\textbf{Отчёт: StackRec — Progressive Layer Stacking для SASRec}\\
\large Эксперименты на датасете Beauty (reviews\_Beauty\_5)}
\author{Проект SASRec experiments}
\date{Октябрь 2026}

\begin{document}

\maketitle

\section*{1. Введение}

StackRec — метод прогрессивного наращивания глубины нейронной сети,
адаптированный для модели SASRec (Sequential Recommendation).
Идея: вместо обучения модели с нуля на полной глубине, обучение
начинается с 2 нижних трансформер-слоёв, затем веса копируются
в новые слои и постепенно активируются все 6 слоёв.

Ключевое свойство: при scale=0 для нижних слоёв они работают как
identity-функция, поэтому подсеть из 2 слоёв в 6-слойной модели
даёт \textbf{абсолютно идентичные} предсказания с самостоятельной
2-слойной SASRec (тесты 16--17, 17/17 PASS).

\section*{2. Архитектура StackRec}

Модель всегда имеет 6 трансформер-слоёв. Per-layer forward coefficients
(\texttt{layer\_scales}) переключают слои: коэффициент 0 делает слой
identity (пропуск), 1 — полная активация.

\begin{itemize}
  \item \textbf{Stage 1 (2\_layers, 90 эпох):} обучаются слои 4,5;
        слои 0--3 заморожены, scale=0.
  \item \textbf{Stage 2 (4\_layers, 60 эпох, 10-epoch ramp):}
        копируем веса 4→2, 5→3; обучаем слои 2--5;
        scale слоёв 2,3 растёт от 0 до 1.
  \item \textbf{Stage 3 (6\_layers, 60 эпох, 10-epoch ramp):}
        копируем 2→0, 3→1; обучаем все 6 слоёв;
        scale слоёв 0,1 растёт от 0 до 1.
\end{itemize}

Каждый stage использует меньший LR по формуле $lr(L) = lr_{ref} \cdot \sqrt{4/L}$.

\section*{3. Эксперименты}

\subsection*{3.1. Базовые SASRec (AdamW / Muon, L=2,4,6)}

\begin{table}[H]
\centering
\begin{tabular}{lccc}
\toprule
Optimizer & L=2 & L=4 & L=6 \\
\midrule
AdamW & 0.012762 & 0.004797 & 0.005645 \\
Muon & 0.005187 & 0.003540 & 0.003692 \\
\bottomrule
\end{tabular}
\caption{Test NDCG@10 для базового SASRec.}
\end{table}

Увеличение глубины ухудшает качество для обоих оптимизаторов.
AdamW стабильно лучше Muon.

\subsection*{3.2. AdamW + StackRec}

\begin{table}[H]
\centering
\begin{tabular}{lcccccc}
\toprule
Layers & LR & Epoch & Val NDCG@10 & Test NDCG@10 & Val Rec@10 & Test Rec@10 \\
\midrule
2 (4,5) & 0.00297 & 88 & 0.014714 & 0.011958 & 0.029630 & 0.023380 \\
4 (2--5) & 0.00210 & 129 & 0.015451 & 0.012964 & 0.031746 & 0.025370 \\
6 (all) & 0.00171 & 194 & 0.016095 & 0.013299 & 0.033439 & 0.025494 \\
\bottomrule
\end{tabular}
\caption{Результаты AdamW + StackRec.}
\end{table}

\begin{table}[H]
\centering
\begin{tabular}{lcc}
\toprule
Layers & AdamW select & StackRec \\
\midrule
2 & 0.012762 & 0.011958 \\
4 & 0.004797 & 0.012964 \\
6 & 0.005645 & 0.013299 \\
\bottomrule
\end{tabular}
\caption{Сравнение: AdamW без stacking vs AdamW + StackRec.}
\end{table}

StackRec даёт в 2.3--2.7 раза лучший test NDCG@10 для 4 и 6 слоёв.

\subsection*{3.3. Muon + StackRec}

\begin{table}[H]
\centering
\begin{tabular}{lcccccc}
\toprule
Layers & Muon LR & AdamW LR & Epoch & Val NDCG@10 & Test NDCG@10 \\
\midrule
2 & 0.02828 & 0.00297 & 85 & 0.007215 & 0.004602 \\
4 & 0.02000 & 0.00210 & 119 & 0.007884 & 0.004518 \\
6 & 0.01633 & 0.00171 & 173 & 0.007008 & 0.004341 \\
\bottomrule
\end{tabular}
\caption{Результаты Muon + StackRec.}
\end{table}

Muon + StackRec улучшает plain Muon для 4 и 6 слоёв, но остаётся
значительно ниже AdamW + StackRec. Лучший Muon stage — 4 layers.

\subsection*{3.4. Auxiliary Identity Loss}

При добавлении новой группы слоёв добавлен auxiliary loss, который
приводит отношение выхода новой группы к входному представлению
обученной подсети к 1:
\[
\lambda(\alpha) = \lambda_0 \cdot (1 - \alpha),
\]
где $\alpha$ — forward scale новой группы.

\begin{table}[H]
\centering
\small
\begin{tabular}{lccc}
\toprule
Opt. & Layers & Val NDCG@10 (plain/aux) & Test NDCG@10 (plain/aux) \\
\midrule
AdamW & 2 & 0.014714 / 0.014714 (+0.00\%) & 0.011958 / 0.011958 (+0.00\%) \\
AdamW & 4 & 0.015451 / 0.015920 (+3.04\%) & 0.012964 / 0.012934 (-0.23\%) \\
AdamW & 6 & 0.016095 / 0.016328 (+1.45\%) & 0.013299 / 0.013086 (-1.60\%) \\
Muon & 2 & 0.007215 / 0.007215 (+0.00\%) & 0.004602 / 0.004602 (+0.00\%) \\
Muon & 4 & 0.007884 / 0.007327 (-7.06\%) & 0.004518 / 0.004591 (+1.61\%) \\
Muon & 6 & 0.007008 / 0.006883 (-1.79\%) & 0.004341 / 0.004351 (+0.24\%) \\
\bottomrule
\end{tabular}
\caption{Влияние auxiliary identity loss.}
\end{table}

Для AdamW: улучшение val NDCG@10 (+3.0\%, +1.5\%), но небольшое ухудшение test.
Для Muon: mixed/negative результаты.

\subsection*{3.5. Hyperbolic Heads}

Poincar\'e-ball next-item head на выходе каждой новой группы:
\[
c(d) = c_{ref} \cdot (d_{ref} / d)^p,
\]
где $d$ — абсолютная глубина head. Кривизна убывает с глубиной.

\begin{table}[H]
\centering
\begin{tabular}{lccc}
\toprule
Layers & Plain & Aux & Hyperbolic \\
\midrule
2 & 0.011958 & 0.011958 & 0.011958 \\
4 & 0.012964 & 0.012934 & 0.012428 \\
6 & 0.013299 & 0.013086 & 0.013419 \\
\bottomrule
\end{tabular}
\caption{Test NDCG@10: plain vs aux vs hyperbolic.}
\end{table}

6-layer модель получает прирост +0.9\% vs plain, +2.5\% vs aux.
4-layer stage слабее — возможно, loss\_coeff=0.1 слишком мал.

\subsection*{3.6. Lorentz Hyperbolic Residual}

Гиперболическая геометрия \textit{в скрытом представлении} копируемой группы:
\[
h_{out} = \text{normalize}(h + \text{gate} \cdot \text{normalize\_to\_norm}(h, r)),
\]
где $r = (\text{sinhc}(\sqrt{c}\|u\|) - 1) \cdot W_{out} u$ — curvature-induced
correction, ровно 0 при $c=0$. Gate и curvature ramp от 0 до целевых значений.

Верификация: unit-level checks (CPU) — hyperboloid constraint, symmetric
distance, zero residual at c=0, bitwise identity at gate=0, norm preservation.

\section*{4. Тестирование}

\subsection*{4.1. Тесты метрик ранжирования (10/10 PASS)}

\begin{table}[H]
\centering
\begin{tabular}{lcl}
\toprule
Метрика & Статус & Значение \\
\midrule
NDCG@5 & PASS & 0.487914 \\
Recall@5 & PASS & 0.800000 \\
Coverage@5 & PASS & 1.000000 \\
Novelty@5 & PASS & 0.944260 \\
MeanRank@5 & PASS & 3.500000 \\
ILD@5 & PASS & 0.506127 \\
PopularityGroups & PASS & OK \\
EdgeCases & PASS & OK \\
CoveragePartial & PASS & 0.100000 \\
MeanRankPartial & PASS & 1.000000 \\
\bottomrule
\end{tabular}
\caption{Тесты метрик ранжирования.}
\end{table}

\subsection*{4.2. Проверка temporal split датасета}

\begin{table}[H]
\centering
\begin{tabular}{lccc}
\toprule
Split & Rows & Users & Items \\
\midrule
train & 140 965 & 18 474 & 11 534 \\
validation & 50 336 & 4 725 & 10 310 \\
test & 90 621 & 8 041 & 11 727 \\
train\_validation & 151 081 & 18 474 & 11 736 \\
\bottomrule
\end{tabular}
\caption{Распределение датасета Beauty 5-core.}
\end{table}

Нарушений temporal split: 0. Prefix-свойство $train \subset val \subset test$: True.

\subsection*{4.3. Тесты StackRec (17/17 PASS)}

\begin{table}[H]
\centering
\small
\begin{tabular}{cl}
\toprule
Тест & Статус \\
\midrule
Model initialization & PASS \\
Forward pass (short seq=5) & PASS \\
Forward pass (long seq=20) & PASS \\
Stage 1 — 2 layers (4, 5) & PASS \\
Stage 2 — 4 layers (2--5) with copy & PASS \\
Stage 3 — 6 layers (0--5) with copy & PASS \\
Full StackRec pipeline (short seq=5) & PASS \\
Full StackRec pipeline (long seq=20) & PASS \\
Ramp scale=0 identity & PASS \\
Aux scalar and boundary & PASS \\
Lorentz head addition & PASS \\
Curvature schedule & PASS \\
Copy layer weights & PASS \\
Batch with two sequences & PASS \\
CPU only & PASS \\
SASRec vs StackRec (2 layers) identity & PASS \\
SASRec vs StackRec (2 layers) long seq & PASS \\
\bottomrule
\end{tabular}
\caption{Тесты StackRec (17 тестов, CPU).}
\end{table}

\textbf{Ключевые проверки:}
\begin{itemize}
  \item \textit{Ramp scale=0 identity:} при scale=0 новые слои — identity,
        предсказания совпадают с моделью, обученной только на 2 слоях.
  \item \textit{SASRec vs StackRec identity:} StackRec с 6 слоями (scale=0
        для 0--3) даёт абсолютно идентичные предсказания с SASRec на 2 слоях.
  \item \textit{Copy layer weights:} веса копируются корректно,
        trainability переключается правильно.
\end{itemize}

\section*{5. Выводы}

\begin{enumerate}
  \item \textbf{StackRec эффективен:} AdamW + StackRec даёт в 2.3--2.7 раза
        лучший test NDCG@10 для 4 и 6 слоёв по сравнению с plain AdamW.
  \item \textbf{Muon + StackRec:} улучшает plain Muon, но остаётся ниже AdamW.
  \item \textbf{Auxiliary identity loss:} улучшает validation для AdamW,
        но test NDCG немного хуже. Для Muon — mixed/negative.
  \item \textbf{Hyperbolic heads:} depth-decaying curvature помогает на
        6 layers (+0.9\% vs plain), но 4-layer stage слабее.
  \item \textbf{Lorentz residual:} верифицирован на unit-тестах,
        full training не запущен.
  \item \textbf{Тестирование:} 10/10 метрик PASS, 17/17 StackRec тестов PASS.
\end{enumerate}

\end{document}